% !TITLE 琵琶行
% !AUTHOR 白居易
% !DYNASTY 唐
% !GENRE 七言古诗
% !ORDER 47

\begin{poem}
浔阳江\textsuperscript{1}头夜送客，枫叶荻花秋瑟瑟\textsuperscript{2}。
主人下马客在船，举酒欲饮无管弦\textsuperscript{3}。
醉不成欢惨将别，别时茫茫江浸月。
忽闻水上琵琶声，主人忘归客不发。
寻声暗问弹者谁？琵琶声停欲语迟\textsuperscript{4}。
移船相近邀相见，添酒回灯重开宴\textsuperscript{5}。
千呼万唤始出来，犹抱琵琶半遮面\textsuperscript{6}。
转轴拨弦\textsuperscript{7}三两声，未成曲调先有情。
弦弦掩抑声声思\textsuperscript{8}，似诉平生不得志。
低眉信手续续弹\textsuperscript{9}，说尽心中无限事。
轻拢慢捻抹复挑\textsuperscript{10}，初为霓裳后六幺\textsuperscript{11}。
大弦嘈嘈如急雨\textsuperscript{12}，小弦切切如私语\textsuperscript{13}。
嘈嘈切切错杂弹，大珠小珠落玉盘\textsuperscript{14}。
间关莺语花底滑\textsuperscript{15}，幽咽泉流冰下难\textsuperscript{16}。
冰泉冷涩弦凝绝\textsuperscript{17}，凝绝不通声暂歇。
别有幽愁暗恨生，此时无声胜有声\textsuperscript{18}。
银瓶乍破水浆迸\textsuperscript{19}，铁骑突出刀枪鸣\textsuperscript{20}。
曲终收拨当心画\textsuperscript{21}，四弦一声如裂帛\textsuperscript{22}。
东船西舫悄无言\textsuperscript{23}，唯见江心秋月白。
沉吟放拨插弦中，整顿衣裳起敛容\textsuperscript{24}。
自言本是京城女，家在虾蟆陵\textsuperscript{25}下住。
十三学得琵琶成，名属教坊\textsuperscript{26}第一部。
曲罢曾教善才服\textsuperscript{27}，妆成每被秋娘妒\textsuperscript{28}。
五陵年少\textsuperscript{29}争缠头\textsuperscript{30}，一曲红绡不知数\textsuperscript{31}。
钿头银篦击节碎\textsuperscript{32}，血色罗裙翻酒污\textsuperscript{33}。
今年欢笑复明年，秋月春风等闲度\textsuperscript{34}。
弟走从军阿姨死，暮去朝来颜色故\textsuperscript{35}。
门前冷落鞍马稀，老大嫁作商人妇\textsuperscript{36}。
商人重利轻别离，前月浮梁\textsuperscript{37}买茶去。
去来江口守空船，绕船月明江水寒。
夜深忽梦少年事，梦啼妆泪红阑干\textsuperscript{38}。
我闻琵琶已叹息，又闻此语重唧唧\textsuperscript{39}。
同是天涯沦落人，相逢何必曾相识\textsuperscript{40}！
我从去年辞帝京，谪居\textsuperscript{41}卧病浔阳城。
浔阳地僻无音乐，终岁不闻丝竹声\textsuperscript{42}。
住近湓江\textsuperscript{43}地低湿，黄芦苦竹绕宅生。
其间旦暮闻何物？杜鹃啼血猿哀鸣。
春江花朝秋月夜，往往取酒还独倾\textsuperscript{44}。
岂无山歌与村笛？呕哑嘲哳难为听\textsuperscript{45}。
今夜闻君琵琶语，如听仙乐耳暂明。
莫辞更坐弹一曲，为君翻作琵琶行\textsuperscript{46}。
感我此言良久立，却坐促弦弦转急\textsuperscript{47}。
凄凄不似向前声，满座重闻皆掩泣\textsuperscript{48}。
座中泣下谁最多？江州司马青衫湿\textsuperscript{49}。
\end{poem}

\begin{annotations}
\notetext{1}{浔阳江：长江流经九江的一段。浔阳（xún yáng），今江西九江。}
\notetext{2}{瑟瑟：秋风的声音，也形容秋风的寒意。枫叶和荻花在秋风中瑟瑟作响。荻（dí），一种水生植物。}
\notetext{3}{无管弦：没有音乐助兴。管弦，管乐和弦乐，代指音乐。}
\notetext{4}{欲语迟：想要开口说话却迟疑了。欲语，想说话；迟，犹豫。}
\notetext{5}{回灯重开宴：把灯重新拨亮，再次摆开宴席。回灯，挑灯使更亮。}
\notetext{6}{犹抱琵琶半遮面：还抱着琵琶遮住半边脸。形容琵琶女初次露面时的羞涩和矜持，后来成了成语。}
\notetext{7}{转轴拨弦：拧动琴轴（调音），拨动琴弦。她先调了调音，还没开始弹，情感已经流露出来。}
\notetext{8}{掩抑声声思：低沉压抑，每一声都带着深沉的思绪。掩抑，压抑不舒展。}
\notetext{9}{信手续续弹：随手连续不断地弹奏。信手，随手；续续，连续不断。}
\notetext{10}{轻拢慢捻抹复挑：四种琵琶弹奏指法。拢，左手扣弦；捻（niǎn），揉弦；抹，右手下拨；挑，右手上挑。}
\notetext{11}{霓裳后六幺：《霓裳羽衣曲》和《六幺》，唐代两大名曲。霓裳是杨贵妃常跳的舞曲，六幺（liù yāo）也是当时流行的大曲。}
\notetext{12}{大弦嘈嘈如急雨：粗弦声音粗重急促，像暴风骤雨打窗。大弦，琵琶最粗的弦。}
\notetext{13}{小弦切切如私语：细弦声音轻细柔和，像人在私下说悄悄话。小弦，琵琶最细的弦。}
\notetext{14}{大珠小珠落玉盘：大大小小的珍珠落在玉盘上——比喻琵琶声清脆明亮如珠玉跳跃。}
\notetext{15}{间关莺语花底滑：琵琶声宛转流利，像黄莺在花丛底下啼叫。间关，鸟鸣声。}
\notetext{16}{幽咽泉流冰下难：琵琶声低沉缓慢，像泉水在冰层下艰难地流动。幽咽（yè），若有若无的低微声音。}
\notetext{17}{冰泉冷涩弦凝绝：像冰封的泉水一样冷涩，琴弦似乎凝固断了。}
\notetext{18}{此时无声胜有声：音乐中断了，但空白处的哀怨比声音更有感染力。这是最有名的音乐描写名句。}
\notetext{19}{银瓶乍破水浆迸：像银瓶突然破裂，水浆飞溅——琵琶声突然爆发。迸（bèng），喷射。}
\notetext{20}{铁骑突出刀枪鸣：像铁甲骑兵突然冲出，刀枪交击声震天——琵琶声的第二轮爆发。}
\notetext{21}{曲终收拨当心画：曲子弹完，收拢拨子，在琴弦中心用力一划。拨，弹琵琶的拨片。}
\notetext{22}{四弦一声如裂帛：四根弦同时发声，像撕裂丝绸一般清脆响亮。}
\notetext{23}{东船西舫悄无言：左右的船只都悄然无声（被音乐震撼了）。舫（fǎng），船。}
\notetext{24}{敛容：收敛笑容，表情变得庄重严肃。弹琴前是羞涩的，弹完后神情变得认真。}
\notetext{25}{虾蟆陵：长安城东南的一个地名，相传是汉代学者董仲舒的墓地。虾（há），同“蛤”。}
\notetext{26}{教坊：唐代官方管理音乐、舞蹈的机构。第一部，第一班，最高级别的演出队。}
\notetext{27}{善才服：连“善才”（曲师、高手）都佩服。善才，琵琶大师的名字，后泛指技艺高超的乐师。}
\notetext{28}{秋娘妒：被“秋娘”（美女）嫉妒。秋娘是唐代歌妓的代称。她的美貌让同行都羡慕嫉妒。}
\notetext{29}{五陵年少：指长安的富家子弟。五陵是汉代五座帝陵（长陵、安陵等），陵区附近多住贵族。}
\notetext{30}{缠头：古代歌舞演员表演后接受的赏赐，常用锦缎缠在头上，故称“缠头”。}
\notetext{31}{一曲红绡不知数：弹完一曲收到的红绡（赏赐的红色丝织品）多得数不清。绡（xiāo），生丝薄纱。}
\notetext{32}{钿头银篦击节碎：用镶嵌珠宝的发饰随着节拍敲击，太投入了，发饰敲碎了。钿（diàn），金花；篦（bì），密齿梳。}
\notetext{33}{血色罗裙翻酒污：红色的罗裙在狂欢中被泼洒的酒弄脏了。写年轻时灯红酒绿的生活。}
\notetext{34}{秋月春风等闲度：美好的时光随随便便地消磨掉了。等闲，随便、轻易。}
\notetext{35}{颜色故：容颜衰老了。故，陈旧、不再新鲜。弟弟参军走了，养母去世了，自己的美貌也一天天褪去。}
\notetext{36}{老大嫁作商人妇：年纪大了，只得嫁给商人做妻子。唐代商人地位低，嫁商人比嫁文人或官吏差很多。}
\notetext{37}{浮梁：今江西景德镇附近，唐代以产茶闻名。丈夫去浮梁买茶，把她一个人留在江边的空船上。}
\notetext{38}{红阑干：胭脂色（泪水）纵横流淌。阑干，纵横交错。梦中哭醒，眼泪冲刷了脂粉。}
\notetext{39}{重唧唧：再次感叹不已。唧唧（jī），叹息声。}
\notetext{40}{同是天涯沦落人：都是流落天涯的人。白居易从京城被贬到浔阳，琵琶女从京城名妓沦落为商人妇——命运相似，所以产生了共鸣。}
\notetext{41}{谪居：被贬谪（zhé）后居住在（浔阳）。白居易因上书言事触怒权贵，被赶出长安。}
\notetext{42}{丝竹声：弦乐器（丝）和管乐器（竹）的声音，代指音乐。终年听不到像样的音乐。}
\notetext{43}{湓江：又名湓水，流经九江的一条小河。住处低洼潮湿，黄芦和苦竹长满了院子。}
\notetext{44}{独倾：独自倒酒独饮。常常在花朝月夜独自一人喝酒，与琵琶女独守空船遥相呼应。}
\notetext{45}{呕哑嘲哳难为听：呕哑（ōu yā）指山歌粗哑，嘲哳（zhāo zhā）指笛声嘈杂刺耳。都难听得没法听。}
\notetext{46}{翻作琵琶行：按照曲调写成这篇《琵琶行》的诗。翻，根据曲调填词。}
\notetext{47}{却坐促弦弦转急：重新坐下，把弦拧紧，弹得更急更快。促，拧紧。}
\notetext{48}{掩泣：掩面哭泣。满座的人再次听到琵琶声，都忍不住哭了。}
\notetext{49}{江州司马青衫湿：江州司马（白居易自称）的青衫（低级官员的官服）被泪水打湿得最厉害。青衫，唐代八、九品官服为青色。}
\end{annotations}

\begin{appreciation}
浔（xún）阳江就是长江流经江西九江的一段。荻（dí）是一种水生植物，瑟（sè）瑟形容秋风吹动的声音。白居易写这首诗时，正因得罪权贵被贬为江州司马，是一个有名无实的闲官。那个秋夜，他送客到江边，偶然听到琵琶声，于是有了这场改变中国文学史的相遇。

诗的前半部分写音乐。大弦嘈嘈如急雨，小弦切切如私语——粗弦声音厚重，像急雨敲窗；细弦声音轻细，像私下低语。白居易用了一连串比喻来写琵琶声：大珠小珠落玉盘，间关莺语花底滑，幽咽泉流冰下难，银瓶乍破水浆迸，铁骑突出刀枪鸣。这些比喻不是随便选的，它们对应着音乐的节奏变化：从流畅到艰涩，从沉寂到爆发。最有名的是此时无声胜有声——音乐停了，但情感还在流动，空白比声音更有力量。

诗的后半部分写琵琶女的身世。她本是京城名妓，十三岁学成琵琶，名属教坊第一部。教坊是唐代管理宫廷音乐的机构。钿（diàn）头银篦（bì）是镶嵌珠宝的发饰。年轻时，她一曲弹完，五陵年少争着送缠头（赏赐），一曲红绡不知数。但好景不长，弟走从军阿姨死，暮去朝来颜色故——弟弟参军走了，养母去世了，年华老去，容颜衰退。最后老大嫁作商人妇，商人重利轻别离，前月浮梁买茶去——丈夫去浮梁买茶，把她一个人留在江口空船上。

我闻琵琶已叹息，又闻此语重唧唧——白居易听了琵琶已经叹息，听了她的身世更加感慨。同是天涯沦落人，相逢何必曾相识——这是全诗的灵魂。一个曾经的京城名妓，一个被贬的朝廷官员，身份悬殊，命运却相似：都是从繁华跌落，都是孤独漂泊。白居易没有居高临下地同情，而是把自己放进去，说“同是”。这种平等的姿态，让这首诗超越了普通的悯人诗，成为两个沦落人之间的相互取暖。

座中泣下谁最多？江州司马青衫湿——青衫是唐代低级官员的服色。白居易哭的不只是琵琶女，也是他自己。那个秋夜的江州司马，在琵琶声中听到了自己的命运。多年后他回首往事，或许会想起那个夜晚：枫叶荻花，江月茫茫，一个弹琵琶的女子，一个听琵琶的诗人，在命运的河流中短暂相遇，然后各自远去。但那琴声和泪水，却留在了诗里，千年不散。
\end{appreciation}
